\documentclass[a4paper,12pt]{article}
\usepackage[T2A]{fontenc}
\usepackage[utf8]{inputenc}
\usepackage[russian]{babel}
\usepackage{amsmath,amssymb,mathtools}
\usepackage{helvet}
\renewcommand{\familydefault}{\sfdefault}
\usepackage{icomma}
\usepackage[left=23mm,right=22mm,top=23mm,bottom=23mm]{geometry}
\usepackage{microtype}
\usepackage{tikz}
\usetikzlibrary{arrows.meta,positioning,shapes.geometric}
\usepackage{tabularx,booktabs,array}
\usepackage{enumitem}
\usepackage{fancyhdr}
\usepackage{setspace}
\usepackage{xcolor}
\usepackage{hyperref}
\definecolor{accent}{HTML}{008080}
\definecolor{darkaccent}{HTML}{004D4D}
\definecolor{textgray}{HTML}{333333}
\hypersetup{colorlinks=true,urlcolor=darkaccent,linkcolor=darkaccent}
\setlength{\parindent}{0pt}
\setlength{\parskip}{5.5pt}
\onehalfspacing
\setlist[itemize]{leftmargin=5.5mm,itemsep=2pt,topsep=2pt}
\setlist[enumerate]{leftmargin=7mm,itemsep=7pt,topsep=4pt}
\pagestyle{fancy}
\fancyhf{}
\fancyhead[L]{\footnotesize\color{textgray} Финансовые задачи \textbullet{} Ксения Клыкова}
\fancyhead[R]{\footnotesize\color{textgray} 02.11.2026}
\fancyfoot[C]{\footnotesize\thepage}
\renewcommand{\headrulewidth}{0pt}
\renewcommand{\footrulewidth}{0pt}
\setlength{\headheight}{15pt}
\newcommand{\mysection}[1]{\vspace{4pt}{\large\bfseries\color{darkaccent}#1}\par\vspace{1pt}}
\newcommand{\idea}[1]{\textcolor{darkaccent}{\textbf{#1}}}
\newcommand{\R}{\,\text{руб.}}
\begin{document}
\begin{titlepage}
\thispagestyle{empty}
\centering
\vspace*{20mm}
{\color{darkaccent}\large МАТЕМАТИКА\par}
\vspace{25mm}
{\LARGE Чек-лист\par}
\vspace{13mm}
{\color{accent}\LARGE Дифференцированная схема\par}
{\color{accent}\LARGE погашения кредита\par}
\vspace{12mm}
{\large Финансовые задачи профильного ЕГЭ\par}
\vspace{25mm}
{\large Ксения Клыкова\par}
\vspace{12mm}
{2 ноября 2026 г.\par}
\vspace{15mm}
{\small Лёвин Артём Александрович\par}
{\small эксклюзивно для Ксении Клыковой\par}
\vspace{5mm}
\begin{tikzpicture}
\draw[accent,line width=0.85pt] (-5.9,0) .. controls (-4,0.8) and (-2,-0.8) .. (0,0)
.. controls (2,0.8) and (4,-0.8) .. (5.9,0);
\end{tikzpicture}
\vspace{9mm}
\end{titlepage}

\mysection{1. Как распознать финансовую схему}
Одно из важнейших действий в финансовой задаче --- правильно прочитать, \emph{какая величина остаётся постоянной}.

\begin{tabularx}{\linewidth}{@{}p{49mm}X@{}}
\toprule
\textbf{Схема} & \textbf{Сигнал в условии} \\
\midrule
Вклад & Деньги размещают в банке; проценты увеличивают накопления; возможны ежегодные пополнения.\\
Аннуитет & Кредит погашается \textbf{равными платежами}: вся выплачиваемая сумма каждый период одинакова.\\
Дифференцированная & Основной долг погашается \textbf{равными долями}; платежи отличаются из-за процентов на уменьшающийся остаток.\\
\bottomrule
\end{tabularx}

\idea{Главное различие:} равная доля долга и равный платёж --- разные условия. При дифференцированной схеме постоянной остаётся доля основного долга, а итоговый платёж убывает.

\mysection{2. Обозначения и единицы}
\begin{tabularx}{\linewidth}{@{}p{20mm}X@{}}
\toprule
$S$ & первоначальная сумма кредита (тело кредита); \\
$n$ & число расчётных периодов и платежей; \\
$r$ & процентная ставка \emph{за один расчётный период}, в процентах; \\
$p=\dfrac{r}{100}$ & та же ставка в виде десятичной дроби; \\
$A_k$ & платёж за $k$-й период; \\
$T$ & сумма всех платежей; \\
$Q=T-S$ & переплата в денежных единицах.\\
\bottomrule
\end{tabularx}

Если платежи ежемесячные, то $n$ --- количество \textbf{месяцев}, а $p$ --- ставка \textbf{за месяц}. Если платежи ежегодные, используем количество лет и годовую ставку. Ставка должна соответствовать периодичности начисления процентов.

\clearpage
\mysection{3. Механика дифференцированного кредита}
Банк выдаёт $S$ на $n$ периодов. В каждый период основная сумма долга уменьшается на одинаковую величину:
\[
\boxed{\text{Доля основного долга}=\frac{S}{n}.}
\]
В начале $k$-го периода остаётся основной долг
\[
S-\frac{(k-1)S}{n}=\frac{S(n-k+1)}{n}.
\]
На него начисляются проценты. Поэтому \textbf{каждый платёж} состоит из двух частей:
\[
\boxed{A_k=\underbrace{\frac{S}{n}}_{\text{доля долга}}+
\underbrace{\frac{S(n-k+1)}{n}p}_{\text{проценты за период}}.}
\]

\mysection{4. Таблица для трёх лет}
Пусть $S=120\R$, $n=3$, $r=10\%$ в год; тогда $p=0,1$, доля долга равна $40\R$.

\renewcommand{\arraystretch}{1.28}
\begin{tabularx}{\linewidth}{@{}c r r r r r@{}}
\toprule
Год & Долг в начале & Проценты & Долг с \% & Платёж & Остаток \\
\midrule
1 & $120$ & $12$ & $132$ & $40+12=52$ & $80$ \\
2 & $80$ & $8$ & $88$ & $40+8=48$ & $40$ \\
3 & $40$ & $4$ & $44$ & $40+4=44$ & $0$ \\
\bottomrule
\end{tabularx}

Сумма выплат: $52+48+44=144\R$. Переплата равна $144-120=24\R$.

\idea{Самопроверка:} в конце срока остаток основного долга равен нулю; суммы частей долга дают $40+40+40=120$; платежи уменьшаются на $4\R$ за год.

\clearpage
\mysection{5. Пять основных формул}
Все формулы ниже справедливы при \textbf{постоянной ставке} $p$ за период и равных частях основного долга.
\begin{align*}
\text{Все платежи:}\quad &\boxed{T=S+\frac{Sp(n+1)}{2}},\\[5pt]
\text{Переплата:}\quad &\boxed{Q=\frac{Sp(n+1)}{2}},\\[5pt]
\text{Доля переплаты:}\quad &\boxed{q=\frac{Q}{S}=\frac{p(n+1)}{2}},\\[5pt]
\text{Первый платёж:}\quad &\boxed{A_1=\frac{S}{n}+Sp},\\[5pt]
\text{Последний платёж:}\quad &\boxed{A_n=\frac{S}{n}+\frac{Sp}{n}}.
\end{align*}

\mysection{6. Откуда берётся формула переплаты}
Проценты начисляются на последовательно уменьшающийся основной долг:
\[
Sp,\quad \frac{S(n-1)}{n}p,\quad \frac{S(n-2)}{n}p,\quad\ldots,\quad\frac{S}{n}p.
\]
Выносим $Sp/n$ за скобки:
\[
Q=\frac{Sp}{n}\bigl(n+(n-1)+\ldots+1\bigr).
\]
Сумма натуральных чисел от $1$ до $n$ равна $n(n+1)/2$, поэтому
\[
Q=\frac{Sp}{n}\cdot\frac{n(n+1)}{2}=\boxed{\frac{Sp(n+1)}{2}}.
\]
Так как все части основного долга в сумме составляют $S$, получаем $T=S+Q$.

\mysection{7. Пример: сумма платежей и процент переплаты}
Кредит $150\R$ выдан на три года под $20\%$ годовых с погашением основного долга равными частями.
\[
S=150,\quad n=3,\quad p=0,2,\quad \frac{S}{n}=50.
\]
Первый платёж: $A_1=50+150\cdot0,2=80\R$. Затем $70\R$ и $60\R$.
\[
T=150+\frac{150\cdot0,2\cdot4}{2}=210\R,
\quad Q=210-150=60\R.
\]
Доля переплаты: $q=60/150=0,4$. \textbf{Процент переплаты: $40\%$.}

\mysection{8. Пример: найти ставку из процента переплаты}
За $14$ месяцев основной долг погашают равными частями. Переплата составляет $15\%$ первоначального кредита. Найдите \emph{месячную} процентную ставку.

\[
q=\frac{15}{100}=0,15,\qquad n=14.
\]
Используем $q=p(n+1)/2$:
\[
0,15=\frac{15p}{2}\quad\Longrightarrow\quad
p=\frac{0,15\cdot2}{15}=0,02.
\]
Следовательно, \textbf{ставка равна $2\%$ в месяц}.

\idea{Важно:} $0,15\cdot2/15=0,02$. Ошибка в положении десятичной запятой меняет ответ в десять раз. Годовую ставку из этой задачи без дополнительного условия вычислять не требуется.

\mysection{9. Почему платежи убывают}
За каждый период основной долг сокращается на $S/n$, поэтому начисленные проценты каждый раз уменьшаются на одну и ту же величину:
\[
\boxed{A_k-A_{k+1}=\frac{Sp}{n}>0\quad\text{при }p>0.}
\]
Это удобный способ проверить полученную последовательность платежей.

\mysection{10. Для сравнения: аннуитетная схема}
При аннуитетном кредите постоянен \textbf{полный платёж} $X$, состоящий из процентов и погашения части долга. Пусть $k=1+p$. Для $n$ равных платежей в конце каждого периода:
\[
\boxed{X=\frac{Sp(1+p)^n}{(1+p)^n-1}},\qquad p>0.
\]
Если $p=0$, то $X=S/n$. Общая сумма платежей $T=nX$, переплата $Q=nX-S$.

\textbf{Пример.} $S=210\R$, два года, $10\%$ годовых, равные ежегодные платежи:
\[
X=\frac{210\cdot0,1\cdot1,1^2}{1,1^2-1}=121\R.
\]
Проверка по остатку: $210\cdot1,1-121=110$, затем $110\cdot1,1-121=0$. Значит, переплата равна $2\cdot121-210=32\R$.

\mysection{11. Для сравнения: вклад с пополнениями}
Пусть на счёте в начале $B_0$, ставка $p$ за период, а \emph{в конце} каждого периода вносится одна и та же сумма $C$. Тогда
\[
\boxed{B_n=B_0(1+p)^n+C\frac{(1+p)^n-1}{p}},\quad p>0.
\]
Если $B_0=0$, $p=0,1$, $n=2$ и требуется накопить $231\R$, то
\[
231=C\cdot\frac{1,1^2-1}{0,1}=2,1C
\quad\Longrightarrow\quad C=110\R.
\]
Проверка: после первого пополнения $110\R$, после второго: $110\cdot1,1+110=231\R$.

\mysection{12. Как избежать типичных ошибок}
\begin{itemize}
\item \textbf{«Равные доли»} не означает «равные платежи». Сначала определите, какая часть суммы фиксирована.
\item \textbf{Проценты начисляются на остаток} основного долга, а не каждый раз на первоначальную сумму.
\item \textbf{Периоды согласованы:} 14 ежемесячных платежей --- это $n=14$; ставка $p$ в формуле также месячная.
\item \textbf{Различайте} $p=0,15$ и $r=15\%$: первое --- десятичная дробь, второе --- процентная запись той же величины.
\item \textbf{Большие коэффициенты:} сначала умножьте уравнение на $10$, $100$ и т.~д., чтобы убрать десятичные дроби, затем сокращайте общие множители.
\item \textbf{Квадратное уравнение:} при вычислении ставки или коэффициента роста оставляйте только корень, соответствующий условиям задачи (в частности, $1+p>0$).
\end{itemize}

\clearpage
\thispagestyle{plain}
\begin{center}
{\LARGE\color{darkaccent}\bfseries Как выбрать финансовую схему}\par
\vspace{5mm}
{\small Сначала определяем постоянную величину, затем выбираем формулу.}
\end{center}
\vspace{8mm}
\begin{center}
\begin{tikzpicture}[>=Latex,
start/.style={ellipse,draw=darkaccent,thin,minimum width=40mm,minimum height=10mm,align=center,fill=white},
question/.style={diamond,draw=darkaccent,thin,aspect=2.7,inner sep=2pt,text width=34mm,align=center,font=\small},
result/.style={rounded corners=3pt,draw=accent,thin,minimum width=46mm,minimum height=12mm,align=center,text width=41mm,font=\small},
line/.style={-{Latex[length=2.5mm]},thin,color=darkaccent},
every node/.style={font=\small}]
\node[start] (begin) at (-2.1,0) {Читаем условие};
\node[question] (q1) at (-2.1,-2.3) {Деньги копятся на счёте?};
\node[result] (r1) at (3.8,-2.3) {Вклад\\Начисление процентов и пополнения};
\node[question] (q2) at (-2.1,-5) {Равны полные платежи?};
\node[result] (r2) at (3.8,-5) {Аннуитет\\Один платёж $X$};
\node[question] (q3) at (-2.1,-7.7) {Равны доли основного долга?};
\node[result] (r3) at (3.8,-7.7) {Дифференцированная схема\\Доля $S/n$ постоянна};
\node[start] (unclear) at (-2.1,-10.3) {Уточнить условия};
\draw[line] (begin) -- (q1);
\draw[line] (q1.east) -- node[above]{да} (r1.west);
\draw[line] (q1.south) -- node[right]{нет} (q2.north);
\draw[line] (q2.east) -- node[above]{да} (r2.west);
\draw[line] (q2.south) -- node[right]{нет} (q3.north);
\draw[line] (q3.east) -- node[above]{да} (r3.west);
\draw[line] (q3.south) -- node[right]{нет} (unclear.north);
\end{tikzpicture}
\end{center}
\clearpage
\mysection{Тренировочный блок --- Путь А}
\textit{10 заданий, от простых к более сложным. Во всех задачах ставка постоянна; начисление процентов и платежи происходят в конце каждого указанного периода, если не сказано иное. В расчётах можно использовать калькулятор.}
\vspace{4mm}
\begin{enumerate}
\item Кредит $120$ тыс. руб. выдан на $3$ года под $10\%$ годовых. Основной долг погашают равными долями. Найдите три ежегодных платежа.
\item Кредит $150$ тыс. руб. выдан на $3$ года под $20\%$ годовых с дифференцированными платежами. Найдите общую сумму выплат.
\item Кредит $400$ тыс. руб. выдан на $4$ года под $10\%$ годовых. Основной долг погашают равными частями. Какова переплата в тысячах рублей?
\item При дифференцированном погашении кредита за $5$ лет ставка составляет $8\%$ годовых. Какой процент от первоначальной суммы составляет переплата?
\item Кредит $280$ тыс. руб. выдан на $7$ лет под $5\%$ годовых с погашением долга равными частями. Найдите первый и последний ежегодные платежи.
\item Кредит погашается равными долями ежемесячно в течение $14$ месяцев. Общая переплата составила $15\%$ от суммы кредита. Найдите месячную ставку.
\item За три года заёмщик выплатил $648$ тыс. руб. по кредиту $600$ тыс. руб. Основной долг гасился равными частями. Найдите годовую процентную ставку.
\item Дифференцированный кредит погашается за четыре года. Первый платёж составляет $84$ тыс. руб., последний --- $66$ тыс. руб. Найдите первоначальную сумму кредита и годовую ставку.
\item Кредит $210$ тыс. руб. взят под $10\%$ годовых на два года. Платежи в конце каждого года одинаковы. Найдите ежегодный платёж и общую переплату.
\item На вкладе первоначально $0$ руб. В конце каждого из двух лет вносят одинаковую сумму. Ставка $10\%$ годовых; проценты начисляют до очередного пополнения. Какую сумму в тысячах рублей надо вносить ежегодно, чтобы после второго пополнения получить $231$ тыс. руб.?
\end{enumerate}

\clearpage
\mysection{Ответы к тренировочному блоку}
\vspace{6mm}
\renewcommand{\arraystretch}{1.6}
\begin{tabularx}{\linewidth}{@{}c X@{}}
\toprule
\textbf{№} & \textbf{Ответ} \\
\midrule
1 & $52$, $48$, $44$ тыс. руб. \\
2 & $210$ тыс. руб. \\
3 & $100$ тыс. руб. \\
4 & $24\%$ \\
5 & $54$ тыс. руб. и $42$ тыс. руб. \\
6 & $2\%$ в месяц \\
7 & $4\%$ годовых \\
8 & $240$ тыс. руб.; $10\%$ годовых \\
9 & $121$ тыс. руб.; переплата $32$ тыс. руб. \\
10 & $110$ тыс. руб. \\
\bottomrule
\end{tabularx}
\vspace{14mm}
\mysection{Итоговая памятка}
\begin{tabularx}{\linewidth}{@{}p{48mm}X@{}}
\toprule
\textbf{Условие} & \textbf{Как действовать} \\
\midrule
Равные доли долга & Используйте $S/n$ и проценты на остаток. \\
Равные платежи & Используйте формулу аннуитета. \\
Нужно найти переплату & Вычислите $Q=Sp(n+1)/2$. \\
Нужно найти ставку & Из $q=p(n+1)/2$ выразите $p$. \\
Нужен первый/последний платёж & Примените $A_1=S/n+Sp$, $A_n=S/n+Sp/n$. \\
Вклад с пополнениями & Учитывайте, в начале или конце периода вносят деньги. \\
\bottomrule
\end{tabularx}
\vfill
\begin{center}
{\small\color{textgray}Лёвин Артём Александрович эксклюзивно для Ксении Клыковой}
\end{center}
\end{document}
